% =====================================================================
%  sections/sec2.tex —— 2 Argyres--Douglas 理论的一般性质（原文 pp.3–5）
% =====================================================================
\section{Argyres--Douglas 理论的一般性质}
\label{sec:general}

Argyres--Douglas 理论最早是作为纯 $\SU(3)$ 规范理论库仑分支上某个点的红外理论被发现的 \cite{6}。在该点上，彼此非局域的无质量双荷子同时存在，于是不存在使所有无质量粒子只带电荷的对偶框架，拉氏量描述也就无从谈起。进一步还可以依据超共形代数论证，这个理论必定是相互作用的 SCFT \cite{11}。

人们又在 $\SU(2)$ QCD 的库仑分支上，通过调节质量参数与库仑分支参数找到了更多例子 \cite{11}。例如，只有一个味的 $\SU(2)$ 给出的 AD 理论，与纯 $\SU(3)$ 理论给出的那个是同一个。这个 AD 理论的 Seiberg--Witten 曲线是
\begin{equation}
x^2 = z^3 + m z + u .
\label{eq:2.1}\tag{2.1}
\end{equation}
各种算子的标度维数，由要求 Seiberg--Witten 微分 $\lambda = x\,\dd z$ 的维数为一这一点定出：
\begin{equation}
[x] + [z] = 1 .
\label{eq:2.2}\tag{2.2}
\end{equation}
我们还要求 Seiberg--Witten 曲线中每一项都有相同的维数，于是上述理论中坐标 $x$ 与 $z$ 的标度维数为 $[x]=\tfrac{3}{5}$、$[z]=\tfrac{2}{5}$。由此容易读出理论中各参数的标度维数：$[m]=\tfrac{4}{5}$、$[u]=\tfrac{6}{5}$。$u$ 是相关算子，$m$ 是这个形变对应的耦合常数。「$[m]+[u]=2$」这一点很重要，唯有如此，把上面的变形解释成保持 $\sN{=}2$ 超对称的形变才是自洽的。事实上，对 AD 理论的任一相关算子 $u_i$，谱中都必须有一个耦合常数 $m_i$，使 $[u_i]+[m_i]=2$ \cite{11}。算子的 $U(1)_R$ 荷与其在超共形点处的标度维数由下式相联系：
\begin{equation}
R(u_i) = 2D(u_i) .
\label{eq:2.3}\tag{2.3}
\end{equation}
不过，这个红外的 $U(1)_R$ 与 UV 的 $U(1)_R$ 对称性毫无关系。

值得注意的是：上面这些例子里找到的 AD 理论，参数个数与 UV 理论相同，但它们的标度维数因强量子效应而剧变。在某些情形下，连库仑分支的维数都会改变。例如，上面那个 AD 理论同样可以从纯 $\SU(3)$ 规范理论得到，而 UV 理论的库仑分支维数为二，红外理论却只有一维。类似地，AD 理论的味对称性也可以与 UV 理论很不一样：上面那个 AD 理论没有味对称性，而作为其来源的带一个味的 $\SU(2)$ 模型有一个 $U(1)$ 味对称性。

下面列举文献中已找到的一些 AD 理论。人们在 $\SU(N)$ 与 $\SO(N)$ QCD 的库仑分支上通过调节参数找到了多得多的例子 \cite{12,13}，并且发现了有趣的 ADE 分类。由 $\SU(2)$ QCD 得到的 AD 理论，按奇异纤维的分类被标为 $A_0$、$A_1$、$A_2$ 理论；它们的高秩推广则由 F 理论给出——把多个 D3 膜放到相应类型的奇点上 \cite{14,15,16}。近期，文献 \cite{17} 用二维/四维对应找到了一大类基于一对 Dynkin 图的理论。用 F 理论研究 AD 理论的近期工作还可参看 \cite{18}。

下面介绍与 $\sN{=}2$ SCFT 有关的几个重要量。四维超共形场论有两个重要的中心荷，它们刻画迹反常 \cite{19}：
\begin{equation}
\langle T^\mu_{\ \mu}\rangle = \frac{c}{16\pi^2}\,(\text{Weyl})^2 - \frac{a}{16\pi^2}\,(\text{Euler}),
\label{eq:2.4}\tag{2.4}
\end{equation}
其中
\begin{equation}
(\text{Weyl})^2 = R^2_{\mu\nu\rho\sigma} - 2R^2_{\mu\nu} + \frac{1}{3}R^2,\qquad
(\text{Euler}) = R^2_{\mu\nu\rho\sigma} - 4R^2_{\mu\nu} + R^2 .
\label{eq:2.5}\tag{2.5}
\end{equation}
对弱耦合的 $\sN{=}2$ 规范理论，中心荷可以用矢量多重态的个数 $n_v$ 与超多重态的个数 $n_h$ 表出 \cite{20}：
\begin{equation}
c = \frac{2n_v + n_h}{12},\qquad
a = \frac{5n_v + n_h}{24},\qquad
a - c = \frac{1}{24}\,(n_v - n_h).
\label{eq:2.6}\tag{2.6}
\end{equation}
这里 $a-c$ 正比于「Higgs」分支的维数，在该分支上所有规范对称性都被完全破缺。另有一个有用的普遍公式，把中心荷与算子谱的标度维数联系起来：
\begin{equation}
4(2a - c) = \sum_{i=1}^{r}\bigl(2D(u_i) - 1\bigr).
\label{eq:2.7}\tag{2.7}
\end{equation}

强耦合理论的中心荷通常极难计算。借助拓扑规范理论，文献 \cite{20} 的作者得到了下面的公式：
\begin{equation}
a = \frac{1}{4}R(A) + \frac{1}{6}R(B) + \frac{5}{24}r + \frac{1}{24}h,\qquad
c = \frac{1}{3}R(B) + \frac{1}{6}r + \frac{1}{12}h,
\label{eq:2.8}\tag{2.8}
\end{equation}
其中 $r$ 是库仑分支一般点处矢量多重态的个数，$h$ 是一般点处自由超多重态的个数；$R(A)$ 与 $R(B)$ 是拓扑规范理论中某些测度因子的 R 荷。$R(A)$ 可由库仑分支算子的标度维数求出：
\begin{equation}
R(A) = \sum_{i=1}^{r}\bigl(D(u_i) - 1\bigr),
\label{eq:2.9}\tag{2.9}
\end{equation}
而 $R(B)$ 与 Seiberg--Witten 曲线的判别式有关，一般极难确定。中心荷也可以用超引力对偶来算 \cite{21}。

\zhnote{[译者按（无外部依据）：原文在 (2.6) 之后给「Higgs」加了引号。之所以加引号，是因为这里指的是「规范对称性被完全破缺的那一支」，对纯规范理论它并不是通常意义上的希格斯分支。译文保留引号。]}
